The prevalent byte pair encoding (BPE) model and the unigram model use only the characters in a text to build a vocabulary. The cheapest kinds of grammatical information are information about the part-of-speech of each word, and lemmas of each word, which can be provided by a standard part-of-speech tagger. A controlled evaluation compares these two models with tokenizers that use part-of-speech and lemma information in five languages: English, Spanish, Danish, Tamil, and Turkish. The same Wikipedia text and \nVocabBudget{} type budget are used for all tokenizers. The first version of a BPE tokenizer uses the probability of the part-of-speech transition at each word boundary in the scoring function for merge operations. The second version uses part-of-speech weighting of merges within a content word (noun, verb, or adjective) so that content words stay whole. Finally the Lemma2Word tokenizer uses the lemma to determine where the boundary between stem and affixes should be.

We trained each model on three different samples of text, evaluated the model on its ability to find morpheme boundaries (the gold standard is the lemma-derived boundaries of MorphScore), held-out compression, stability of the model, and bits per character of a small language model trained on the model's output. The part-of-speech weighting of transition probabilities has no effect, in four out of five languages, on the F1 score for morpheme boundary detection (the difference is less than the variation from seed to seed), and of course does not affect held-out compression. This is not surprising, since these affect tokens that are common function word pairs, and of course a part-of-speech transition probability does not convey information about the interior of a word. However the part-of-speech weighting of the merger of tokens within content words raises the F1 score for finding morpheme boundaries from \nFBpeTamil\% to \nFPosContentTamil\% in Tamil, while lowering the scores in Spanish and Turkish. Lemma2Word raises boundary recall in every language by 23 to 77 points (with 4--8\% more tokens). When the tokenizer does not see the word in its training text, the model retains most of the gain if the language's affixes are long enough for an inventory of them to generalize.

The probe with a language model ranks the tokenizers differently. For example, the unigram model has the lowest number of bits per character in every language, while Lemma2Word has the highest. This shows that a tokenizer that finds more morpheme boundaries is not necessarily a better choice for a language model, at least at the scale tested here.
